We can prove the following result using Green\textquotesingle s functions and the theory of compact operators.

\begin{Theorem}
\label{thm:regular-sturm-liouville-complete-spectral}
Let $(\mathcal{D}(L),L)$ be defined by \eqref{eq:sturm-liouville-weighted-differential-operator} and \eqref{eq:sturm-liouville-unmixed-operator-domain}. The following hold:\begin{enumerate}\item The spectrum consists entirely of eigenvalues.\item The eigenvalues are countable and can be listed in a sequence
\begin{equation}\label{eq:sturm-liouville-increasing-eigenvalue-sequence}
\lambda_1<\lambda_2<\cdots<\lambda_n<\cdots
\end{equation}
with
\begin{equation}\label{eq:sturm-liouville-eigenvalues-tend-to-infinity}
\lim_{n\to\infty}\lambda_n=\infty.
\end{equation}
\item The set of normalized eigenfunctions $\{\phi_i\}$ such that $\phi_i\in\mathcal{N}(L_{\lambda_i})$ (where $L_\lambda:=L-\lambda I$) is an orthonormal basis for $L_w^2(a,b)$.\item For the equation
\begin{equation}\label{eq:sturm-liouville-spectral-inhomogeneous-equation}
Lu-\lambda u=f,
\end{equation}
exactly one of the following alternatives holds:\begin{enumerate}[label=(\alph*)]\item If $\lambda$ is not an eigenvalue of $(\mathcal{D}(L),L)$, then \eqref{eq:sturm-liouville-spectral-inhomogeneous-equation} has a unique solution in $\mathcal{D}(L)$ for every $f\in L_w^2(a,b)$. This solution is given by
\begin{equation}\label{eq:sturm-liouville-noneigenvalue-solution-expansion}
u(x)=\sum_{i=1}^\infty\dfrac{(\phi_i,f)_w}{\lambda_i-\lambda}\phi_i(x).
\end{equation}
\item If $\lambda=\lambda_j$ is an eigenvalue of $(\mathcal{D}(L),L)$, then \eqref{eq:sturm-liouville-spectral-inhomogeneous-equation} has a solution in $\mathcal{D}(L)$ provided
\begin{equation}\label{eq:sturm-liouville-eigenvalue-solvability-condition}
(\phi_j,f)=0.
\end{equation}
In this case there is a one-parameter family of solutions given by
\begin{equation}\label{eq:sturm-liouville-eigenvalue-solution-family}
u(x)=C\phi_j(x)+\sum_{\substack{i=1\\i\ne j}}^\infty\dfrac{(\phi_i,f)_w}{\lambda_i-\lambda}\phi_i(x).
\end{equation}
\end{enumerate}\end{enumerate}
\end{Theorem}

\begin{Proof}
We begin by constructing a Green\textquotesingle s function for the ODE
\begin{equation}\label{eq:sturm-liouville-homogeneous-expanded-ode}
-pu''-p'u'+(q-\lambda w)u=0
\end{equation}
and the boundary conditions \eqref{eq:sturm-liouville-left-unmixed-boundary-condition} and \eqref{eq:sturm-liouville-right-unmixed-boundary-condition}. Let $v_l(x;\lambda)$ satisfy the ODE \eqref{eq:sturm-liouville-homogeneous-expanded-ode} and the left-hand boundary condition \eqref{eq:sturm-liouville-left-unmixed-boundary-condition} and $v_r(x;\lambda)$ satisfy \eqref{eq:sturm-liouville-homogeneous-expanded-ode} and the right-hand end condition. We assume that $\lambda$ is not an eigenvalue. In this case, the ODE uniqueness theorem for initial-value problems implies that $v_l$ and $v_r$ are linearly independent.

The Green\textquotesingle s functions will have the form
\begin{equation}\label{eq:sturm-liouville-green-function-piecewise-ansatz}
g(x,y;\lambda):=\begin{cases}
a_l(y;\lambda)v_l(x;\lambda),&a<x<y,\\
a_r(y;\lambda)v_r(x;\lambda),&y<x<b.
\end{cases}
\end{equation}
Thus $g(\cdot,y;\lambda)$ satisfies \eqref{eq:sturm-liouville-left-unmixed-boundary-condition} and \eqref{eq:sturm-liouville-right-unmixed-boundary-condition}. To get
\begin{equation}\label{eq:sturm-liouville-green-function-delta-equation}
-p\frac{\partial^2}{\partial x^2}g(x,y;\lambda)-p'\frac{\partial}{\partial x}g(x,y;\lambda)+(q-\lambda w)g(x,y;\lambda)=\delta(x-y)
\end{equation}
we require\[\begin{aligned}
a_l(y;\lambda)v_l(y;\lambda)-a_r(y;\lambda)v_r(y;\lambda)&=0,\\
a_l(y;\lambda)v_l'(y;\lambda)-a_r(y;\lambda)v_r'(y;\lambda)&=\dfrac{1}{p(y)}.
\end{aligned}\]Solving this gives us
\begin{equation}\label{eq:sturm-liouville-green-left-coefficient}
a_l(y;\lambda)=(W(v_l(y;\lambda),v_r(y;\lambda))p(y))^{-1}v_r(y;\lambda),
\end{equation}

\begin{equation}\label{eq:sturm-liouville-green-right-coefficient}
a_r(y;\lambda)=(W(v_l(y;\lambda),v_r(y;\lambda))p(y))^{-1}v_l(y;\lambda),
\end{equation}
where
\begin{equation}\label{eq:sturm-liouville-fundamental-wronskian}
W(v_l(y;\lambda),v_r(y;\lambda)):=v_l(y;\lambda)v_r'(y;\lambda)-v_l'(y;\lambda)v_r(y;\lambda)
\end{equation}
is the Wronskian of $v_l$ and $v_r$. A classical ODE result called \emph{Abel\textquotesingle s formula} (which we do not prove) gives us
\begin{equation}\label{eq:sturm-liouville-abel-constant-wronskian}
W(v_l(y;\lambda),v_r(y;\lambda))p(y)=C^{-1}
\end{equation}
where $C$ is a constant. Thus, we have
\begin{equation}\label{eq:sturm-liouville-green-function-explicit-kernel}
g(x,y;\lambda):=\begin{cases}
Cv_r(y;\lambda)v_l(x;\lambda),&a<x<y,\\
Cv_l(y;\lambda)v_r(x;\lambda),&y<x<b.
\end{cases}
\end{equation}
Note that $g$ is bounded, and since $w(x)$ is bounded, the kernel $w(y)g(x,y;\lambda)$ is Hilbert-Schmidt. Thus, the operator
\begin{equation}\label{eq:sturm-liouville-weighted-green-resolvent-operator}
L_w^2(a,b)\ni u\mapsto K_\lambda u:=\int_a^bw(y)g(\cdot,y;\lambda)u(y)\,\mathrm{d}y\in L_w^2(a,b)
\end{equation}
is bounded. A direct computation shows that $K_\lambda=R_\lambda(L)$. Since $K_\lambda$ has dense domain, and is well defined if and only if $\lambda$ is not an eigenvalue of $(\mathcal{D}(L),L)$, we have established 1.

By part 2 of the previous theorem we can assume without loss of generality that $\lambda=0$ is not an eigenvalue (otherwise make an appropriate change in $q$). Let $g_0(x,y):=g(x,y;0)$. Note that the integral equation
\begin{equation}\label{eq:sturm-liouville-kernel-eigenproblem}
u(x)=\lambda\int_a^bw(y)g_0(x,y)u(y)\,\mathrm{d}y
\end{equation}
has a nontrivial solution $u$ if and only if $\lambda$ is an eigenvalue of $(\mathcal{D}(L),L)$ and $u$ a corresponding eigenfunction. By letting\[\begin{aligned}
v(x)&:=\sqrt{w(x)}u(x),\\
k(x,y)&:=\sqrt{w(x)}g_0(x,y)\sqrt{w(y)},
\end{aligned}\]we can write \eqref{eq:sturm-liouville-kernel-eigenproblem} in the form
\begin{equation}\label{eq:sturm-liouville-unweighted-compact-eigenproblem}
Gv(x):=\int_a^bk(x,y)v(y)\,\mathrm{d}y=\dfrac{1}{\lambda}v(x).
\end{equation}
Since $k$ is symmetric and bounded, $G$ is a self-adjoint, compact operator from $L^2(a,b)$ to $L^2(a,b)$. Eigenvalues of $(\mathcal{D}(L),L)$ are clearly reciprocals of the eigenvalues of $G$, and normalized eigenfunctions of $L$ can be obtained from those of $G$ by the formula $\phi_i(x)=w(x)^{-1/2}v_i(x)$. Furthermore, since for any $f\in L_w^2(a,b)$ we have $\sqrt{w}f\in L^2(a,b)$, so we can write
\begin{equation}\label{eq:sturm-liouville-weighted-data-basis-expansion}
\sqrt{w(x)}f(x)=\sum_{i=1}^\infty(v_i,\sqrt{w}f)v_i(x)=\sqrt{w(x)}\sum_{i=1}^\infty(\phi_i,f)_w\phi_i(x).
\end{equation}
Hence,
\begin{equation}\label{eq:sturm-liouville-data-eigenfunction-expansion}
f(x)=\sum_{i=1}^\infty(\phi_i,f)_w\phi_i(x).
\end{equation}
This proves 2 and 3.

Existence and multiplicity of solutions for the nonhomogeneous equation follows from the Fredholm alternative theorem and the equivalent integral formulations. To derive the expansion formulas, assume first that $\lambda$ is not an eigenvalue. Let
\begin{equation}\label{eq:sturm-liouville-solution-basis-expansion}
u(x)=\sum_{i=1}^\infty(\phi_i,u)_w\phi_i(x)
\end{equation}
and
\begin{equation}\label{eq:sturm-liouville-forcing-basis-expansion}
f(x)=\sum_{i=1}^\infty(\phi_i,f)_w\phi_i(x).
\end{equation}
Then
\begin{equation}\label{eq:sturm-liouville-equating-spectral-coefficients}
(L-\lambda I)u=\sum_{i=1}^\infty(\phi_i,u)_w(\lambda_i-\lambda)\phi_i(x)=\sum_{i=1}^\infty(\phi_i,f)_w\phi_i(x).
\end{equation}
Equating coefficients give us \eqref{eq:sturm-liouville-noneigenvalue-solution-expansion}. A similar computation yields \eqref{eq:sturm-liouville-eigenvalue-solution-family}. This completes the proof.
\end{Proof}
